\begin{problem}{TST 2025 - Bài 6}{standard input}{standard output}{10 seconds}{256 megabytes}

Alice rất yêu thích các xâu kí tự. Với ba số nguyên dương $N, L, R$ ($1 \le L \le R \le N$), cô đã tạo ra tập $\Omega$ gồm tất cả các xâu độ dài $N$ là hoán vị của $N$ kí tự Latin in thường đầu tiên, sao cho vị trí của kí tự \texttt{a} trên xâu thuộc đoạn $[L; R]$ (vị trí của các ký tự trong xâu được đánh số từ $1$ đến $N$).

Bob muốn tạo ra một xâu $U$ chứa tất cả các xâu trong tập $\Omega$ và có độ dài càng nhỏ càng tốt. Ở đây, xâu $X$ được gọi là \textbf{chứa} xâu $Y$ nếu như có thể xóa bỏ một số ký tự của xâu $X$ (có thể không xóa ký tự nào) và giữ nguyên thứ tự của các ký tự còn lại để có được xâu $Y$.

Tuy nhiên, Alice đã không cho Bob biết chính xác giá trị của $N,L,R$. Thay vào đó, cô sẽ trả lời một số câu hỏi của anh. Với mỗi lần hỏi của Bob:
\begin{itemize}
   \item Bob sẽ đưa cho Alice một xâu $s$ có độ dài không quá $200$. Xâu này chỉ được chứa các ký tự Latin in thường.
   \item Alice sẽ tính toán và trả lời cho Bob biết số xâu $t\in \Omega$ thỏa mãn xâu $s$ khớp với xâu $t$ chia lấy dư cho $379$.
\end{itemize}

Sau khi hỏi Alice một số câu hỏi, Bob sẽ trả lời xâu $U$ mà anh tìm được. Để thể hiện sự thông minh của mình trước Alice, Bob muốn giá trị $P=Q+|U|$ càng nhỏ càng tốt, trong đó $Q$ là số câu hỏi Bob đã thực hiện và $|U|$ là độ dài xâu $U$ mà Bob trả lời. Hãy lập trình giúp Bob một chương trình tương tác với Alice để thực hiện yêu cầu trên.

\InputFile
\textbf{Chi tiết cài đặt}

Thí sinh cần cài đặt hàm sau:

\begin{lstlisting}
string solve()
\end{lstlisting}

\begin{itemize}
   \item Hàm này có thể được gọi nhiều lần bởi trình chấm.
   \item Hàm này cần trả về xâu $U$ thỏa mãn điều kiện đề bài. Lưu ý là bạn không được biết trước giá trị của $N,L,R$.
\end{itemize}

Trong hàm \texttt{solve()}, thí sinh được phép gọi hàm sau (thí sinh cần khai báo trước thư viện \texttt{smatchlib.h} để có thể sử dụng hàm này):

\begin{lstlisting}
int ask(string s)
\end{lstlisting}

\begin{itemize}
   \item $s$: Xâu mà Bob sẽ đưa cho Alice để thực hiện câu hỏi. Xâu $s$ chỉ được chứa các ký tự Latin in thường và có độ dài không quá $200$.
   \item Hàm sẽ trả về một số nguyên không âm là số lượng xâu $t\in \Omega$ thỏa mãn $s$ khớp với $t$ chia lấy dư cho $379$.
\end{itemize}

Ngoài ra, thí sinh được phép cài đặt các biến, các hàm toàn cục, nhưng không được phép có hàm \texttt{main()}.

Cách cài đặt và trình chấm ví dụ có thể xem tại contest material, lưu ý trình chấm ví dụ chỉ mô phỏng cách hoạt động và sẽ không được dùng để chấm chính thức.

\textbf{Ràng buộc}

\begin{itemize}
   \item $4\le N\le 13$
   \item $1\le L\le R\le N$
\end{itemize}

Trình chấm của bài này \textbf{không mang tính thích ứng}. Điều này tương đương với việc các giá trị $N,L,R$ sẽ được cố định từ đầu và không thay đổi trong suốt quá trình chương trình của bạn tương tác với trình chấm.

\Scoring
\begin{center}
  \begin{tabular}{|c|c|l|} \hline
    \textbf{Subtask} & \textbf{Điểm} & \textbf{Giới hạn} \\ \hline
    1 & $10$ & $N=4$ \\ \hline
    2 & $10$ & $N\le 5$ \\ \hline
    3 & $10$ & $N\le 6$ \\ \hline
    4 & $10$ & $N\le 7$ \\ \hline
    5 & $10$ & $N\le 8$ \\ \hline
    6 & $10$ & $N\le 9$ \\ \hline
    7 & $10$ & $N\le 10$ \\ \hline
    8 & $30$ & Không có ràng buộc gì thêm \\ \hline
  \end{tabular}
\end{center}

\textbf{Cách tính điểm}

Giả sử test có tối đa $1$ điểm. Nếu trong bất kỳ lần gọi hàm \texttt{solve} nào, xâu $U$ bạn đưa ra không thỏa mãn với điều kiện của đề bài thì bạn sẽ được $0$ điểm cho toàn bộ test đó. Ngược lại, điểm của bạn sẽ được tính như sau:
\begin{itemize}
   \item Mỗi lần gọi hàm \texttt{solve} sẽ có điểm của lần gọi hàm đó (gọi là $score$). Xét lần gọi hàm tương ứng, gọi giá trị $P=Q+|U|$ là độ tối ưu kết quả của thí sinh, giá trị $J=Q+|U|$ là độ tối ưu kết quả của ban tổ chức ($Q$ là tổng số lần thực hiện truy vấn \texttt{ask}, $U$ là xâu đáp án trả về đã thỏa mãn điều kiện). Khi đó:
   \begin{itemize}
      \item Nếu $Q>10$ thì $score=0$.
      \item Nếu $N<P-J$ thì $score=0.4\times 0.9^{P-J-N-1}$.
      \item Nếu $4\le P-J\le N$ thì $score=0.4+0.4\times 0.8^{P-J-4}$.
      \item Nếu $0<P-J<4$ thì $score=0.8+0.2\times 0.8^{P-J}$.
      \item Nếu $P\le J$ thì $score=1$.
   \end{itemize}
   \item Điểm của một test sẽ bằng điểm (giá trị $score$) nhỏ nhất trong tất cả lần gọi hàm \texttt{solve}.
\end{itemize}

Điểm của mỗi subtask là điểm thấp nhất của một test trong subtask đó.


\end{problem}

